\documentclass[]{article}

\usepackage{listings}
\usepackage{color}

\definecolor{dkgreen}{rgb}{0,0.6,0}
\definecolor{gray}{rgb}{0.5,0.5,0.5}
\definecolor{mauve}{rgb}{0.58,0,0.82}

\lstset{frame=tb,
	language=Matlab,
	aboveskip=3mm,
	belowskip=3mm,
	showstringspaces=false,
	columns=flexible,
	basicstyle={\small\ttfamily},
	numbers=none,
	numberstyle=\tiny\color{gray},
	keywordstyle=\color{blue},
	commentstyle=\color{dkgreen},
	stringstyle=\color{mauve},
	breaklines=true,
	breakatwhitespace=true,
	tabsize=3
}

%opening
\title{Complex Magnitude API}
\author{Kenneth Martin}

\begin{document}

\maketitle

\begin{abstract}
A very concise desription of the functions used.
\end{abstract}

\section*{Functions}
\subsubsection*{cascadeClass.m}
\begin{lstlisting}
    classdef (ConstructOnLoad = true) cascadeClass < handle
    cascadeClass defines a class to hold Cascade Filter instantiations
    The cascadeClass object contains biquads in cscFltrSctnClass objects
    It also has a function for evaluating the frequency response of
    complex cascade filters plus a number of other utility functions.
    sctns: an array of cscFltrSctnClass objects which contain first or second order sections.
    size: the number of sections
    obj = cascadeClass(sctn): make a new cascadeClass object, sctn can be
    empty, or a cascadeClass obj, for copying a filter, or cscFltrSctnClass which gets added as the first section
    obj = addSctn(obj, sctn): add an additional section
    gn = freqEval(obj, f): evaluate the cascade filter as frequencies f (in real Hz)
    gn = scaleFltr(obj, wp): scale gains of sections so gain peaks are all 1 (f-inf scaling)
    sys = getSystem(obj): update and return the overall system in a zpk system obj
    sys = fshftFltr(obj, fshft): freq shift cascade filter; returns shifted zpk obj
    sys = fsclFltr(obj, sclFctr): freq scale cascade filter using scaleFltrD.m
    plotGn(obj, wp, ws, minY, maxY): plot the response of the cascade filter using plot_drsps.m
    out = sim(obj, xin, delta_f): simulate the cascade filter for complex input signal xin, after optionally frequency shifting filter (useful for filterbanks)
\end{lstlisting}

\subsubsection*{cscFltrSctnClass.m}
\begin{lstlisting}
    classdef (ConstructOnLoad = true) cscFltrSctnClass < handle
    The cscFltrSctnClass is a class for a section or biquad of a digital cascade filter
    The sections are first or second order. The sections ar stored as a zpk model and also
    as zeros, poles, and k factor for easy Whenever any sub components are changed,
    the zpk model is updated. There are a number of utility functions such as frequency
    analysis, difference-equation simulation, etc.
    obj = cscFltrSctnClass(sys): make a new cscFltrSctnClass initialized to sys (a zpk
    system obj). If sys is a cscFltrSctnClass object, a copy is returned.
    obj = setSys(obj, sys): change the system to a new zpk sys
    obj = setZ(obj, z): change the zeros, and update sys
    obj = setP(obj, p): change the poles and update sys
    obj = setK(obj, k): change the k factor and update sys
    mxGn = getGain(obj, wp): find the maximum gain aroung wp = [wp1, wp2]
    obj = setGain(obj, gain, wp): set the maximum gain
    gn = freqEval(obj, f): evaluate gain at frequencies f (freqz() is used for evaluation)
    out = sim(obj, xin, delta_f):simulate the filter section for complex input signal xin,
    after optionally frequency shifting filter section by delta_f(useful for filterbanks)
\end{lstlisting}

\subsubsection*{chckEqlOrdr.m}
\begin{lstlisting}
    isEqual = chckEqlOrdr(ply, fun) checks if polyClass object ply is identical to fun. It's used to possibly terminate finding roots of a function early when the fun is itself a polynomial having fewer roots than expected.
\end{lstlisting}

\subsubsection*{chooseX.m}
\begin{lstlisting}
    chooseX(func, inArry) applies min(func) to inArry
    and returns the indx and the array with the min removed
    It's a utility function used in grouping zeros and poles
    when designing a digital filter
\end{lstlisting}

\subsubsection*{cleanRoots.m}
\begin{lstlisting}
    rts2 = cleanRoots(rts, tol) sets any real or imaginary parts of vector rts having an abs() value less than tol to 0. tol defaults to 1e-5. Components with absolute values larger than tol are unchanged.
\end{lstlisting}

\subsubsection*{dAs10\_dz.m}
\begin{lstlisting}
    h = dAs10_dz(wsz,As,w) finds derivate of stopband specs using simple
    interpolation of specifications in log10 to the transformed
    variable z
\end{lstlisting}

\subsubsection*{design\_dig\_filt.m}
\begin{lstlisting}
    [H, E, F, P, e_] = design_dig_filt(p,px,ni,wp,ws,as,ap,type) Design a discrete-time complex IIR transfer function. 
    p: Initial guess at finite loss poles; actual values not important as long as they are in the stop-band; this spec is mostly used to decide the order.
    px: fixed poles; often used to have a loss-pole at dc or the carrier frequency.
    ni: number of poles at infinity of the continuous-time prototype. They result in loss-poles at half the sample frequency
    wp: the predistorted pass-band frequencies in radians. If the final desired passband frequencies (in Hz) are [fp1 fp2], then uses wp = 2*tan(pi.*[fp1 fp2]).
    ws: the stop-band frequencies in radians; at least two, but there can be more than two; the number of frequencies should match the size of as; all loss-pole frequencies (p, and px) should be in the stop-bands
    as: the specifications for desired attenuation at each stop-band frequency; the actual attenuation is determined by the order and how wide the transition region is, but generally the DIFFERENCES in attenaution at the ws frequencies is determined by the as specifications. as values should all be positive
    ap: the desired pass-band ripple in dB; normally this spec. is met exactly.
    type: 'monotonic' for a maximally flat pass-band and 'elliptic' for an equi-ripple pass-band
\end{lstlisting}

\subsubsection*{disp\_margind.m}
\begin{lstlisting}
    [margin, smin] = disp_margind(sys,ws,as,wp,e_,type) prints stop-ban margins, that is the differences between the stopband loss and th interpolated specs at the frequncies specified by the vector w. As should be in dB. The specification frequencies are in the transformed domain.
\end{lstlisting}

\subsubsection*{dlogKK\_dp.m}
\begin{lstlisting}
    h = dlogKK_dp(sys,px,s) is a program for calculating the derivative of the log of the magnitude of system Kz_ with respect to the los poles.
    sys: the LTI object.
    s: the frequencies of the minima
    h: the returned matrix has one row for each minima frequency and one column for each moveable loss-pole
\end{lstlisting}

\subsubsection*{dAs10\_dz.m}
\begin{lstlisting}
    h = dAs10_dz(wsz,As,w) finds derivate of stopband specs using simple
    interpolation of specifications in log10 to the transformed variable z
\end{lstlisting}

\subsubsection*{dsply.m}
\begin{lstlisting}
    dsply(X0) is a utility function used during debug to simplify looking at a system object it is just a cover for sortZPK to be easier to type. z, and p are returned as vectors sorted on the jw axis
\end{lstlisting}

\subsubsection*{dsplyRes.m}
\begin{lstlisting}
    dsplyRes(X0) is a utility function to display the residues of a system object. It is primarily used when debugging ladder removals
\end{lstlisting}

\end{document}
